\scanpage{79}
% archive sheet number: 78
\uline{Scholie} \inserted{Soit $n\in\mathbb N\setminus\{1\}$, d'où $\Lambda_n\simeq\begin{cases}\mathbb Z&\text{si }n=0\\\mathbb Z/n\mathbb Z&\text{sinon}\end{cases}$,} Par les principes généraux, il s'ensuit qu'on a \deleted{\unsure{entre n'importe|réciproquement}} \inserted{\deleted{deux}} \unsure{équivalences|équivalence} de groupoïdes
\[
 (M,i)\longmapsto\widetilde P_{M,i}
\]
\diagram{\begin{tikzcd}[ampersand replacement=\&,column sep=small]
{\text{(18)}} \& {\substack{\unsure{\text{Groupoïde}}\,\inserted{\text{des}}\,\text{modules libres}\\\text{$\frac12$ symplectiques }\deleted{\text{de rang 2}}\inserted{\text{de rang 2}}\text{ sur }\Lambda_n\\\deleted{\unsure{\Lambda_n}}}}\arrow[r,"{\approx}"] \&
{\substack{\text{groupoïde des}\\\text{structures triangulées}\\\text{avec antipodisme}\\\text{(autom involutif)}\\\deleted{\unsure{\text{satisf}}}\text{ isomorphes à}\\\text{une structure du type}\\\text{canonique (dépendant}\\\deleted{\unsure{\text{de }\Lambda}}\text{ de }n\text{)}}}
\end{tikzcd}}
\[
 M\longmapsto\widetilde P_M
\]
\diagram{\begin{tikzcd}[ampersand replacement=\&,column sep=small]
{\text{(19)}} \& {\substack{\text{groupoïde des}\\\text{modules libres de}\\\text{rang 2 sur }\Lambda_n\\\deleted{\unsure{(n=0,\;n\in\mathbb N\setminus\{1\})}}}}\arrow[r,"\approx"] \&
{\substack{\text{groupoïde des}\\\text{structures }(S,F,\unsure{\pi|\Pi},\tau,a)\\\text{où }\deleted{\unsure{(S,F)}}\deleted{\unsure{\text{avec}}}\text{ on a une}\\\deleted{\unsure{F}}F\subset\mathfrak P_3(S)\text{ une str}\\\text{triangulée sur }S,\\\deleted{\unsure{F}}\text{ décomposée en familles}\\\inserted{\text{disjointes}}\text{ indexées par }\pi\\\text{(i.e. on a }\tau:F\longrightarrow\pi\ \text{surj)}\\a\text{ est un autom inv.}\\\text{(antipodisme) de la}\\\text{structure précédente.}}}
\end{tikzcd}}
% slanted
\margin{\uline{NB} Si $2.1_\Lambda=0$ (i.e., ici, $\Lambda=\mathbb F_2$) il est inutile d'introduire l'antipodisme, qui est l'identité. Si $2.1_\Lambda\ne0$ \deleted{\unsure{($\Lambda=\mathbb F_2$ exclu)}} alors l'antipodisme est une involution ss pts fixes...}

Il conviendrait de caractériser axiomatiquement, par des propriétés intrinsèques, les structures qu'on obtient au deuxième mb. De plus, il faudrait comprendre, en termes de ces données, \uline{comment reconstruire} $M$, et par là aussi,
% slanted
\margin{cf suggestion (23) ci-dessous}
